\section{Related Work}
\subsection{Generative Recommendation}

In recent years, generative recommendation has attracted considerable attention in both academia and industry. This emerging paradigm, usually built on the Transformer architecture, reformulates recommendation as an end-to-end next-item generation task and thereby raises the performance ceiling of recommender systems.
Early work TIGER \citep{TIGER} employs residual quantization (RQ-VAE) \citep{rqvae} to convert text embeddings—extracted from an item’s title and description—into discrete semantic IDs, which are then used for next-item prediction. LC-Rec \citep{LCRec} aligns an LLM with semantic IDs via multi-task learning, enabling the model to “understand” the IDs and perform generative recommendation. Other studies investigate how to build better semantic IDs to enhance generation quality: RecForest \citep{recforest} applies hierarchical k-means clustering and treats the cluster indices as tokens, while EAGER \citep{eagar} and TokenRec \citep{TOKENRec} integrate semantic and collaborative signals directly into the tokenizer.

Very recently, generative recommendation has been rolled out at an industrial scale to address the drawbacks of traditional cascade systems. MTGR \citep{mtgr} keeps the original deep learning recommendation model (DLRM) features, introduces user-level compression, and speeds up both training and inference for large-scale deployment. OneRec \citep{OneRec} lowers the serving cost with a Lazy Decoder-Only Architecture and stabilises training through an improved reinforcement learning algorithm.

\subsection{LLM and RL}
Reinforcement Learning (RL) \citep{RLsurvey, RLsurvey2} operates through the interaction mechanism between agents and the environment, rewarding correct actions with the core objective of maximizing cumulative returns. In recent years, researchers have introduced RL into the fine-tuning process of LLMs. Among these approaches, Reinforcement Learning with Human Feedback (RLHF) \citep{RLHF} typically employs the Proximal Policy Optimization (PPO) \citep{PPO} algorithm to achieve alignment with human preferences. However, the PPO algorithm faces technical challenges related to excessive memory consumption. To simplify the application of RL in LLMs, researchers have proposed innovative methods such as Direct Preference Optimization (DPO) \citep{DPO}. In the context of recommender systems, the s-DPO method \citep{sdpo} combines DPO with the softmax loss function, enhancing LLMs' capability in mining hard negative samples for recommender systems. Although these methods have improved efficiency, they still face challenges such as off-policy issues and performance ceilings that fall short of online methods. The recently proposed Group Relative Policy Optimization (GRPO) \citep{deepseekmath} further reduces memory requirements by estimating group scores, while replacing traditional reward models with rule-based reward mechanisms. Through this innovative paradigm, LLMs can significantly enhance their reasoning capabilities and domain-specific performance (e.g., in mathematics and programming) without requiring annotated data \citep{Deepseek-r1, o1}. Recent studies \citep{R1-searcher, rm-r1, search-r1, RecZero} have further explored the applications of RL in mathematics and general domains.

% \subsection{LLM post-training Integration}
% Reinforcement Learning (RL) \citep{RLsurvey, RLsurvey2} operates through the interaction mechanism between agents and the environment, rewarding correct actions with the core objective of maximizing cumulative returns. In recent years, researchers have introduced RL into the fine-tuning process of LLMs. Among these approaches, Reinforcement Learning with Human Feedback (RLHF) \citep{RLHF} typically employs the Proximal Policy Optimization (PPO) \citep{PPO} algorithm to achieve alignment with human preferences. However, the PPO algorithm faces technical challenges related to excessive memory consumption. To simplify the application of RL in LLMs, researchers have proposed innovative methods such as Direct Preference Optimization (DPO) \citep{DPO}. In the context of recommender systems, the s-DPO method \citep{sdpo} combines DPO with the softmax loss function, enhancing LLMs' capability in mining hard negative samples for recommender systems. Although these methods have improved efficiency, they still face challenges such as off-policy issues and performance ceilings that fall short of online methods. The recently proposed Group Relative Policy Optimization (GRPO) \citep{deepseekmath} further reduces memory requirements by estimating group scores, while replacing traditional reward models with rule-based reward mechanisms. Through this innovative paradigm, LLMs can significantly enhance their reasoning capabilities and domain-specific performance (e.g., in mathematics and programming) without requiring annotated data \citep{Deepseek-r1, o1}. Recent studies \citep{R1-searcher, rm-r1, search-r1} have further explored the applications of RL in mathematics and general domains. Despite these advancements, the application of RL in improving rating prediction performance through reasoning capabilities remains unexplored.
